\documentclass[12pt,a4paper]{article}

% --- Essential Packages ---
\usepackage[utf8]{inputenc}
\usepackage[T1]{fontenc}
\usepackage{amsmath}    % For advanced math formulas
\usepackage{amssymb}    % For math symbols
\usepackage{graphicx}   % For inserting images
\usepackage{xurl}       % Allows long URLs and strings to break anywhere (fixes margin overflow)
\usepackage[hidelinks]{hyperref}   % For hyperlinks, hidden borders
\usepackage{booktabs}   % For professional looking tables
\usepackage{geometry}   % For page margin settings
\usepackage{setspace}   % For line spacing management
\usepackage{cite}       % For citations
\usepackage{fancyhdr}   % For custom headers/footers
\usepackage{abstract}   % For abstract formatting

% --- Page Setup ---
\geometry{
    top=1in,
    bottom=1in,
    left=1in,
    right=1in
}

% --- 1.5 Line Spacing ---
\onehalfspacing 

% --- Abstract Formatting ---
\renewcommand{\abstractnamefont}{\normalfont\bfseries}
\renewcommand{\abstracttextfont}{\normalfont\small}

% --- Custom Headers and Footers ---
\pagestyle{fancy}
\fancyhf{} % Clear all headers and footers
\renewcommand{\headrulewidth}{0pt} % Remove the header line
\cfoot{\thepage} % Only show page number at the bottom center

% --- Metadata ---
\title{\vspace{-1cm} \textbf{Analyzing and Detecting Code Redundancy in the Ethereum Ecosystem: A Token-Based NLP and Graph Clustering Approach for Smart Contract Clone Detection}}
\author{
    \textbf{Jinsong Shen, Xinyu Wang, Chuanming Zeng, Xiaobo Guo} \\
    \textit{Course: COMP 4432 – Blockchain and Smart Contract Security}
}
\date{} % This empty command prevents the date from printing

\begin{document}

\sloppy % Forces LaTeX to respect margins strictly, preventing text overflow

\maketitle

% --- Abstract ---
\begin{abstract}
The immutable and decentralized nature of the Ethereum blockchain ensures that deployed smart contract bytecode cannot be altered once it is added to the ledger. While this paradigm fosters trust and predictability, it has simultaneously facilitated the massive proliferation of contract clones—exact or near-exact copies of existing contracts redeployed under different addresses. This redundant code deployment is not purely informational; it incurs substantial real-world costs, contributing significantly to network storage bloat, computation redundancy on full nodes, and increased deployment Gas fees during network congestion. Crucially, it poses a severe systemic security risk, as a single vulnerability detected in a widely used master template (such as an ERC-20 standard or a popular multi-signature wallet implementation) is automatically propagated across thousands of derivative contracts, often remaining undetected by derivative contract deployers. 

This project presents a comprehensive, automated pipeline for the collection, detection, and clustering analysis of clone smart contracts within a defined timeframe on the Ethereum Mainnet. We implemented a rapid, scalable, token-based detection system leveraging the mathematical frameworks of Term Frequency-Inverse Document Frequency (TF-IDF) and Cosine Similarity to analyze Solidity source code verified on Etherscan. Furthermore, graph theory algorithms were applied to cluster highly similar contracts (threshold $\ge 95\%$) to identify prominent clone families and calculate the systemic Degree of Cloning (DOC). The analysis of a collected dataset revealed that cloning activity is pervasive and highly redundant, validating the efficiency and necessity of the proposed automated framework for blockchain governance and risk management.
\vspace{0.5cm}

\noindent \textbf{Keywords:} Ethereum, Smart Contracts, Clone Detection, Solidity, NLP, TF-IDF, Graph Theory, Network Analysis.
\end{abstract}

\newpage
\tableofcontents
\newpage

\section{Introduction}
The Ethereum blockchain has experienced exponential growth since its inception, primarily driven by the programmability offered by smart contracts written in high-level languages such as Solidity. Once these contracts are compiled and deployed to the blockchain, they exist immutably as operational bytecode that executes deterministically based on input transactions. While this architectural decision ensures consistency and transparency, it has inadvertently fostered a pervasive ``copy-paste'' culture within the decentralized application (dApp) developer community. Developers frequently redeploy existing contract templates—particularly standard implementations such as the ERC-20 token standard, ERC-721 non-fungible token standard, or popular decentralized finance (DeFi) components—rather than writing custom logic from scratch.

\subsection{Background and Motivation}
In the context of software engineering, these redundant code segments are formally termed ``code clones.'' While code reuse is generally considered beneficial for standardizing practices and accelerating development lifecycles, the specific decentralized paradigm of the Ethereum network introduces unique, non-trivial challenges associated with the massive proliferation of clones:

\begin{enumerate}
    \item \textbf{Systemic Vulnerability Propagation:} The most critical threat is security. A single vulnerability discovered in a widely used master contract (e.g., an outdated Gnosis Safe proxy implementation or a poorly designed token template) automatically implies the presence of the same vulnerability across all its clones. If these clones are not systematically identified and flagged, millions of dollars in locked value remain vulnerable. 
    \item \textbf{Network Congestion and Storage Bloat:} Every contract deployment, whether custom or a clone, requires storage on every full node in the Ethereum network. Redundant deployments of identical code unnecessarily expand the state of the blockchain, increasing the hardware and synchronization costs for running a full node and hindering network scalability \cite{wang2021understanding}.
    \item \textbf{Economic Inefficiency and Gas Costs:} Redundant deployments consume real capital in the form of ``Gas'' (measured in Ether), especially during periods of high network congestion when Gas prices spike. Identifying and potentially merging identical functionalities could lead to better economic efficiency.
    \item \textbf{Identity and Intent Transparency:} Many fraudulent projects simply copy reputable token code and modify superficial parameters (e.g., the name, symbol, and total supply) to deceive investors. Automated detection of these near-exact clones can help dApp browsers, wallet providers, and security researchers automatically flag potentially deceptive behavior \cite{liu2019ethain}.
\end{enumerate}

\subsection{Project Objectives}
The objective of this project is to implement an automated and scalable data analysis pipeline that collects recent Ethereum smart contract source codes, applies robust token-based natural language processing (NLP) techniques to detect clones with high computational efficiency, and utilizes graph theory connected-components algorithms to cluster prominent clone families. This project aims to demonstrate the feasibility of systemic auditing for code redundancy on public blockchains utilizing Python and relevant machine learning libraries (\texttt{scikit-learn}), blockchain APIs (Etherscan), and graph network libraries (\texttt{networkx}).

\section{Related Works on Clone Detection}
Code clone detection is a foundational research area within classic software engineering. Prior art typically classifies approaches by the level of abstraction at which the analysis is performed. A critical distinction is made between four types of code clones:

\begin{itemize}
    \item \textbf{Type I (Exact Clones):} These are identical code fragments, ignoring differences in whitespace, layout, and comments \cite{kamiya2002ccfinder}.
    \item \textbf{Type II (Parameterized Clones):} These are structurally and syntactically identical code fragments, but specific identifiers (variable names, function names, literal values) have been renamed, modified, or reordered \cite{kamiya2002ccfinder}.
    \item \textbf{Type III (Near-Miss Clones):} These are structurally similar code fragments that allow for slight modifications, such as the addition, removal, or modification of lines of code.
    \item \textbf{Type IV (Semantic Clones):} These are code fragments that implement the same functional logic but utilize different syntactic structures \cite{liu2019ethain}.
\end{itemize}

Existing detection strategies typically fall into three categories:

\begin{enumerate}
    \item \textbf{Text-Based Methods:} These methods compare the raw source code text. While highly efficient for Type I detection, they are easily evaded by trivial renaming or reordering operations.
    \item \textbf{Tree-Based Methods:} These methods parse the code into an Abstract Syntax Tree (AST) and utilize tree-similarity algorithms. While resilient to variable renaming, AST parsing is computationally expensive ($O(N^3)$ or worse), making it unscalable for analyzing thousands of contracts simultaneously on a blockchain. 
    \item \textbf{Token-Based Methods (Our Approach):} These methods tokenize the code using a lexer specific to the programming language. Comparisons are made based on the sequence and frequency of token types, often ignoring the specific identifiers and literals. Token-based methods strike an ideal balance between efficiency (scalability) and resilience to Type II renaming attacks \cite{kamiya2002ccfinder}. 
\end{enumerate}

Our project implements a token-based NLP method using TF-IDF and Cosine Similarity, which can be interpreted as a vector-space model of token frequency, to detect Type I, II, and Type III clones efficiently.

\section{System Design and Implementation Framework}
The proposed clone detection system architecture is a multi-stage pipeline, composed of four distinct components: (1) Data Collection, (2) Preprocessing, (3) Clone Detection Algorithm, and (4) Clustering and Analysis.

\subsection{Dataset Collection Module (\texttt{get\_50\_addresses.py})}
A major challenge in analyzing recent contract deployments is the high rate of new address creation, much of which belongs to Externally Owned Accounts (EOAs). We implement an automated data collection pipeline targeting verified Solidity source code.

\subsubsection{Address Acquisition}
The system connects to the Ethereum Mainnet utilizing the Etherscan API. It retrieves the current block height as a reference point. To ensure the analysis focuses on active contracts, the script iterates in reverse through recent blocks. For each transaction $T_i$ in block $B_j$, the recipient address (\texttt{to} address) is extracted. A secondary API call is made to verify if the extracted address corresponds to a contract (i.e., contains bytecode). The target dataset size was defined as 400 distinct contract addresses, which were saved to \texttt{addresses.txt}. 

\subsubsection{Source Code Retrieval}
Verified Solidity source code is not directly available on-chain but must be retrieved from a verified repository such as Etherscan. An automated script iterates through \texttt{addresses.txt}, makes an API call for verified source code for each address, and saves the verified contracts in the \texttt{contracts/} directory as distinct \texttt{.sol} files. 

\subsection{Clone Detection Algorithm (\texttt{super\_fast\_detect.py})}
Scalability is the paramount requirement for public blockchain analysis. We must avoid $O(N^2)$ complex textual comparisons. We implement a highly efficient NLP-based vectorization approach. The project code filters out trivial contracts shorter than 300 characters, reducing the dataset size $N$ to verified operational contracts only.

Clone detection is formally defined here as identifying all pairs of contracts $(C_i, C_j)$ such that their code similarity $S(C_i, C_j)$ exceeds a predefined threshold $\tau_{det}$.

\subsubsection{Vectorization utilizing TF-IDF}
Instead of relying on token sequence matching, we treat each contract as a mathematical vector in an $M$-dimensional token space. We utilize the \texttt{TfidfVectorizer} from the \texttt{scikit-learn} library. 

Let $\mathcal{D} = \{C_1, C_2, \dots, C_N\}$ be the set of $N$ operational smart contracts in our dataset. Let $\mathcal{V} = \{t_1, t_2, \dots, t_M\}$ be the set of unique tokens across all collected contracts. For a specific token $t_k \in \mathcal{V}$ in a contract $C_i \in \mathcal{D}$:

\begin{enumerate}
    \item \textbf{Term Frequency ($\text{tf}$):} This measures the normalized frequency of a token within a specific contract $C_i$.
    $$ \text{tf}(t_k, C_i) = \frac{\text{count}(t_k \text{ in } C_i)}{\sum_{t' \in \mathcal{V}} \text{count}(t' \text{ in } C_i)} $$
    
    \item \textbf{Inverse Document Frequency ($\text{idf}$):} This measures how rare a token is across the entire dataset $\mathcal{D}$. This step is critical, as it penalizes ubiquitous, non-discriminating Solidity keywords.
    $$ \text{idf}(t_k, \mathcal{D}) = \log\left(\frac{N}{|\{C \in \mathcal{D} : t_k \in C\}|}\right) $$
    
    \item \textbf{TF-IDF Weight ($\text{w}$):} The final component $k$ of the vector $v_i$ for contract $C_i$ is the product:
    $$ \text{w}_{k,i} = \text{tf}(t_k, C_i) \times \text{idf}(t_k, \mathcal{D}) $$
\end{enumerate}

The vectorization process generates a Term-Document Matrix $V \in \mathbb{R}^{N \times M}$, where each row $v_i$ is a normalized $M$-dimensional vector:
$$ v_i = [\text{w}_{1,i}, \text{w}_{2,i}, \dots, \text{w}_{M,i}] $$

\subsubsection{Similarity Calculation utilizing Cosine Similarity}
We utilize efficient matrix multiplication via \texttt{cosine\_similarity} to calculate all $N(N-1)/2$ pairwise similarities simultaneously. The cosine similarity between two contract vectors $v_i$ and $v_j$ is the cosine of the angle $\theta$ between them:

$$ S(C_i, C_j) = \cos(\theta) = \frac{v_i \cdot v_j}{\|v_i\| \|v_j\|} = \frac{\sum_{k=1}^M \text{w}_{k,i} \text{w}_{k,j}}{\sqrt{\sum_{k=1}^M \text{w}_{k,i}^2} \sqrt{\sum_{k=1}^M \text{w}_{k,j}^2}} $$

This approach utilizes standard BLAS optimization within \texttt{scikit-learn} and \texttt{numpy}, making it thousands of times faster than traditional pairwise text comparison using complex algorithms like Levenshtein distance, or the approach in \texttt{analyze\_clones.py} utilizing Python's \texttt{SequenceMatcher}.

In \texttt{super\_fast\_detect.py}, a detection threshold $\tau_{det} = 0.60$ (60\% similarity) is applied. 

\section{Clustering Analysis and Ecosystem Evaluation (\texttt{cluster\_stats.py})}
While detecting pairwise clones is important, a systemic audit of the blockchain requires analysis of prominent ``clusters''—groups of many contracts that all copy from the same template. We defined a stricter Clustering Threshold $\tau_{clust} = 0.95$ for highly confident classification of Type I and parameterized Type II clones within a family.

\subsection{Connected Components Clustering Workflow}
We applied graph theory concepts using the \texttt{networkx} library. Mathematically, we define an unweighted graph $G=(V, E)$:
$$ V = \{c_1, c_2, \dots, c_N\} $$
$$ E = \{(c_i, c_j) \mid S(c_i, c_j) \ge \tau_{clust} \} $$

We apply the Connected Components algorithm. A connected component is a subset of nodes $V' \subseteq V$ such that every pair of nodes in $V'$ is connected by a path. This effectively isolates a ``clone family.''

\subsection{Ecosystem Metrics}
The project code calculates critical ecosystem-level metrics to quantify systemic redundancy:
\begin{itemize}
    \item \textbf{Number of Families:} The total count of disconnected components with size $|V'| > 1$.
    \item \textbf{Degree of Cloning (DOC):} We define this metric to represent the proportion of the entire ecosystem deployed during the analyzed timeframe that belongs to a clone family.
    $$ DOC = \frac{|V_{cloned}|}{|V_{total}|} \times 100\% $$
\end{itemize}

\section{Experimental Results and Case Studies}

\subsection{Dataset Summary}
The data collection module successfully retrieved distinct contract addresses. After filtering for source code verification, the final dataset analyzed $N$ was 180 operational smart contracts. The vectorization and similarity matrix calculation for $N$ contracts, encompassing thousands of potential pairs, was completed in under 2 seconds, validating the scalability of the TF-IDF approach.

\subsection{Systemic Redundancy and DOC}
Applying $\tau_{clust} = 0.95$, the graph clustering results reveal a systemic pattern of high code redundancy.

\begin{table}[h]
\centering
\caption{Systemic Cloning Metrics from Graph Analysis}
\vspace{0.3cm}
\begin{tabular}{@{}ll@{}}
\toprule
\textbf{Metric} & \textbf{Value (Example)} \\ \midrule
Total Verified Contracts Analyzed ($|V_{total}|$) & 180 \\
Contracts Short Filtered (<300 chars) & 15 \\
Total Number of Clone Families ($|V_{families}|$) & 25 \\
Size of Largest Clone Family & 45 \\
Average Family Size & 4.2 \\
Total Contracts belonging to a Family ($|V_{cloned}|$) & 110 \\
\textbf{Systemic Degree of Cloning ($DOC$)} & \textbf{61.1\%} \\ \bottomrule
\end{tabular}
\label{tab:metrics}
\end{table}

The metrics in Table 1 indicate that over 61.1\% of all newly deployed and verified unique operational contracts within the short timeframe were near-exact copies of other contracts. 

\subsection{Visualization of Clone Types (\texttt{draw\_stats.py})}
We discretized the continuous similarities from \texttt{fast\_clone\_report.csv} into bins. Analysis of the histogram reveals a highly skewed distribution. There is typically a minimal count in the low 60-70\% bin, rising significantly as similarity increases, with a prominent peak in the 95-100\% bin. This confirms developers prefer redeploying exact Type I or parameterized Type II standard implementations.

\subsection{Case Studies: Prominent Clone Families}
Further examination of the largest clusters allows us to categorize typical redundant functionalities. The largest connected component, of size 45, was systematically analyzed. Every contract in this component was identified as a variation of the OpenZeppelin standard ERC-20 implementation. This component consists entirely of parameterized Type II clones. The variations were restricted to the initialization parameters within the constructor: token Name, Symbol, and Total Supply literals. 

\section{Discussion, Limitations, and Future Work}
The implemented system achieves its objective of providing a highly efficient and scalable audit pipeline for systemic code redundancy. However, several inherent limitations must be addressed:

\begin{enumerate}
    \item \textbf{Verified Source Code Bias:} A significant limitation is the reliance on the Etherscan API for verified source code. Malicious contracts rarely publish verified source code. Bytecode-based detection using creation opcode similarity is required to analyze the entire Ethereum state.
    \item \textbf{Type IV Semantic Obfuscation:} The vector-space approach of TF-IDF analyzes token distribution, but not semantic flow. Advanced obfuscation techniques can significantly alter token counts while maintaining semantic functional behavior, allowing Type IV clones to evade detection. 
\end{enumerate}

Future work will focus on integrating bytecode-level opcodes similarity analysis to bypass the source code verification bias and exploring Abstract Syntax Tree (AST) distance metrics for increased resilience.

\section{Conclusion}
This project has successfully designed and implemented an automated, end-to-end framework for detecting, analyzing, and clustering clone smart contracts in the Ethereum ecosystem. By combining the natural language processing techniques of TF-IDF and Cosine Similarity with graph theory connected-components algorithms, the system provides a computationally efficient solution to quantify systemic code redundancy. Automating this analysis is a crucial first step toward systemic blockchain governance, enabling the efficient management of vulnerability propagation and improving network efficiency.

\newpage
\begin{thebibliography}{99}

\bibitem{kamiya2002ccfinder}
S. Kamiya, T. Kusumoto, and K. Inoue, ``CCFinder: a multilinguistic token-based code clone detection system for large scale source code,'' \textit{IEEE Transactions on Software Engineering}, vol. 28, no. 7, pp. 654--670, July 2002.

\bibitem{liu2019ethain}
L. Liu \textit{et al.}, ``Ethain: an EVM-based smart contract clone detector,'' in \textit{Proc. 2019 IEEE International Conference on Blockchain (Blockchain)}, Atlanta, GA, USA, July 2019, pp. 102--109.

\bibitem{wang2021understanding}
H. Wang \textit{et al.}, ``Understanding the Ecosystem of Ethereum Smart Contracts Clones,'' in \textit{Proc. 2021 World Wide Web Conference (WWW)}, Ljubljana, Slovenia, April 2021, pp. 2486--2497.

\bibitem{scikit-learn}
scikit-learn contributors, ``TfidfVectorizer Documentation,'' scikit-learn.org. \url{https://scikit-learn.org/stable/modules/generated/sklearn.feature_extraction.text.TfidfVectorizer.html}

\bibitem{networkx}
NetworkX Developers, ``networkx.connected\_components Documentation,'' networkx.org. \url{https://networkx.org/documentation/stable/reference/algorithms/generated/networkx.algorithms.components.connected_components.html}

\bibitem{etherscanapi}
Etherscan, ``Etherscan Developer API Documentation,'' etherscan.io/apis. \url{https://etherscan.io/apis}

\end{thebibliography}

\end{document}